\begin{enumerate}
\item Nếu $(Q)$ chứa $a$ và song song với $(P)$: qua một điểm của $a$ dựng thêm đường thẳng $b$ song song với $(P)$, $b$ cắt $a$; khi đó $(Q)=(a,b)$.
\item Nếu $(Q)$ chứa $a$ (hoặc đi qua hai điểm) và song song với đường thẳng $d$: qua một điểm của $a$ dựng $d'\parallel d$, khi đó $(Q)=(a,d')$.
\item Tìm giao tuyến của $(Q)$ với các mặt của hình, dùng giao tuyến song song, rồi nối các đoạn giao tuyến để được thiết diện.
\end{enumerate}
Lưu ý: thiết diện là đa giác khép kín, mỗi cạnh nằm trên một mặt của hình; mỗi mặt chứa nhiều nhất một cạnh của thiết diện.
